% s-192-p192-m001; original PDF 192, 193
\Aphorism{s-192-p192-m001}{279}
Das Werk! Wir kennen das Werk als Erleichterung, als Arbeit, als umgeformtes Quantum von Zeit und Technik: kennen wir es auch als Selbstzerstörung, als Fatum, als das grosse Einmal-und-nicht-wieder, das seinen Urheber verschwendet, ihm erlaubt, nahelegt, gebietet zu verschwinden? Es giebt ein Schaffen, das an Passivität grenzt; man reimt ein Sonett, wie man etwas aus der Tasche verliert — wir geben nichts von uns hin, es gleitet nur etwas von uns ab. Es giebt ein Schaffen, das schon einen wirklichen Kraftverlust darstellt, das uns in Anspruch nimmt, aber nicht erschöpfend, vielleicht nur unser Gehirn, irgend etwas Einzelnes an uns. Endlich das ganze Schaffen, die Erschöpfung des Schaffenden; nicht mehr verwandelte Energie und partielle Hingabe, die sich ersetzen lässt, sondern ernstgenommene Unmöglichkeit, weiterzuleben, der Mensch ins Werk aufgelöst ohne Rückstand und Vorbehalt. Man hüte sich hier an vorhandene Dinge zu denken; bis jetzt hat noch jeder Schaffende seine Schöpfung überlebt, noch kein Werk den Wirkenden zerschmettert. Wir müssen den natürlichen Leichtsinn verehren, der immer wieder „hinanzieht“, nämlich aus den Tiefen der eigentlichen Production zur Oberfläche des detaillirten Ausarbeitens — und von dort lässt sich allerlei Land gewinnen. Aber schweben uns nicht bisweilen Werke vor, nach denen es uns natürlich schiene, in der persönlichen Existenz eine Pause eintreten zu lassen, Dinge, die wir mit dem gleichen selbstvernichtenden Ernst betreiben könnten wie manche Insekten die Begattung: Geständnisse, die eine Seele dem Leben gegenüber unrettbar compromittiren müssten, Worte, die nur als Formeln eines letzten Willens verzeihlich wären? Dieser extreme Fall, bis heute unerlebt, wäre der einzige, der es gestattete, das Werk seinem Urheber zuzurechnen, wogegen es uns gleichgültig sein kann, zu einem Stück Arbeit den Verfertiger, zu einer verlorenen Sache den Eigenthümer zu wissen.
\par

